\documentclass[../../main.tex]{subfiles}
\begin{document}

{\bf [解答]}

題意より
\begin{align}
  f(x) = \int_0^x \cos t \, dt - 2 \int_0^x \sin t \, dt \quad \left(0 < x < \dfrac{\pi}{4}\right)
\end{align}
とおくと，$f(x)>0$を示せば良い．
一階微分は
\begin{align}
  f'(x) = \cos x - 2 \sin x
\end{align}
であり，これは区間 $[0, \pi/4]$ で単調減少かつ，$f'(\pi/6) = 0$ より
$f(x)$の増減表は以下のようになる．

\begin{center}
  \begin{table}[htb]
    \centering
    \caption{$f(x)$の増減表（$0 < x < \pi/4$）}
    \begin{tabular}{c|ccccc}
      $x$  & $0$ & $\dots$ & $\pi/6$ & $\dots$ & $\pi/4$ \\ \hline
      $f'$ & & $+$ & $0$ & $-$ & \\ \hline
      $f$  & $0$& $\nearrow$ & & $\searrow$ & $\dfrac{3}{\sqrt{2}} - 2$
    \end{tabular}
  \end{table}
\end{center}

これと
\begin{align}
  \begin{cases}
    f(0)=0 \\
    f(\pi/4) = [\sin t + 2 \cos t]_0^{\pi/4} = \dfrac{3}{\sqrt{2}} - 2 > 0
  \end{cases}
\end{align}
より$f(x) > 0 \ (0 < x < \pi/4)$ なので，題意は示された．

\end{document}
